\section{Methods}
\label{sec:methods}

We describe the mathematical formulation and GPU implementation strategy for each of the five computational kernels.
Throughout, we adopt the following notation conventions: $\gamma$ denotes the Lorentzian half-width at half-maximum (HWHM) in \si{\nano\metre}, $\sigma$ denotes the Gaussian standard deviation in \si{\nano\metre}, $T$ denotes the electron temperature in \si{\electronvolt} unless subscripted as $T_K$ for Kelvin, $n_e$ denotes the electron density in \si{\per\cubic\centi\metre}, and $\lambda$ denotes wavelength in \si{\nano\metre}.
Energies $E_k$ are measured in \si{\electronvolt} above the ground state of each ionization stage, and $k_B = \SI{8.617333e-5}{\electronvolt\per\kelvin}$.

\subsection{GPU-Optimized Voigt Profile}
\label{sec:voigt}

The Voigt profile $V(\lambda)$ describes the spectral line shape arising from the convolution of Gaussian (Doppler) and Lorentzian (Stark broadening plus natural linewidth) contributions.
In terms of the Faddeeva function $w(z) = e^{-z^2} \operatorname{erfc}(-iz)$, the normalized Voigt profile is
\begin{equation}
  V(\lambda) = \frac{\operatorname{Re}[w(z)]}{\sigma \sqrt{2\pi}}, \quad
  z = \frac{(\lambda - \lambda_0) + i\gamma}{\sigma\sqrt{2}},
  \label{eq:voigt}
\end{equation}
where $\lambda_0$ is the line center wavelength.
The profile integrates to unity over wavelength and has dimensions of \si{\per\nano\metre}.

Efficient evaluation of $w(z)$ on GPU hardware requires a branch-free algorithm: conditional branches within vectorized operations cause warp divergence on NVIDIA architectures and are incompatible with JAX's automatic differentiation, which evaluates all branches of \texttt{jnp.where} during the forward pass.
The widely used \citet{Humlicek1982} W4~algorithm employs four piecewise regions selected by conditionals on $|x|$ and $y = \operatorname{Im}(z)$, making it unsuitable for JAX compilation.

We adopt the rational approximation of \citet{Weideman1994} with $N = 36$ terms.
This algorithm applies a M\"obius transformation $Z = (L + iz)/(L - iz)$ with optimal shift $L \approx 5.0454$, then evaluates a degree-35 polynomial in $Z$:
\begin{equation}
  w(z) \approx \frac{1}{\sqrt{\pi}} \frac{L}{L - iz}
    + \frac{2}{(L - iz)^2} \sum_{n=0}^{35} c_n Z^n,
  \label{eq:weideman}
\end{equation}
where the coefficients $\{c_n\}$ are precomputed constants.
This formulation is entirely branch-free, requires 236 floating-point operations per evaluation, and achieves a maximum relative error below $10^{-13}$ in 64-bit arithmetic across the entire upper half-plane $\operatorname{Im}(z) \geq 0$ \citep{Weideman1994}.
\citet{Zaghloul2024} recently proposed an improved algorithm that achieves $10^{-16}$ accuracy, which we use as an accuracy reference; however, its additional computational cost is unnecessary for LIBS applications where the Voigt profile is convolved with an instrument function of finite resolving power.

For a spectrum containing $N_\mathrm{lines}$ emission lines evaluated on a wavelength grid of $N_\mathrm{wl}$ points, the full profile matrix has shape $(N_\mathrm{wl}, N_\mathrm{lines})$.
Rather than looping over lines, we broadcast the wavelength grid against the line parameter vectors:
\begin{equation}
  \mathrm{diff}_{ij} = \lambda_i - \lambda_{0,j}, \quad
  z_{ij} = \frac{\mathrm{diff}_{ij} + i\gamma_j}{\sigma_j \sqrt{2}},
  \label{eq:voigt-broadcast}
\end{equation}
where $i = 1, \ldots, N_\mathrm{wl}$ and $j = 1, \ldots, N_\mathrm{lines}$.
The entire $(N_\mathrm{wl}, N_\mathrm{lines})$ grid of Faddeeva evaluations is then executed as a single fused XLA kernel, exploiting the massive parallelism of the GPU.

\subsection{Vectorized Boltzmann Fitting}
\label{sec:boltzmann}

The Boltzmann plot method \citep{Tognoni2010} estimates the plasma temperature from emission line intensities.
For a given element and ionization stage, plotting
\begin{equation}
  y = \ln\!\left(\frac{I \lambda}{g_k A_{ki}}\right)
  \quad \text{vs.} \quad
  x = E_k
  \label{eq:boltzmann-plot}
\end{equation}
yields a linear relationship $y = bx + a$ with slope $b = -1/(k_B T)$ and intercept $a = \ln(hc\, C_s\, n_e\, F_s / (4\pi\, U_z(T)))$, where $I$ is the measured line intensity, $g_k$ is the upper-level degeneracy, $A_{ki}$ is the Einstein~$A$ coefficient, $E_k$ is the upper-level energy, and $F_s$ encodes Saha ionization-stage ratios.

Rather than an iterative regression or matrix factorization, we express the weighted least-squares (WLS) solution in closed form via five dot-product accumulators:
\begin{equation}
  S_w = \sum_i w_i, \quad
  S_{wx} = \sum_i w_i x_i, \quad
  S_{wy} = \sum_i w_i y_i, \quad
  S_{wxx} = \sum_i w_i x_i^2, \quad
  S_{wxy} = \sum_i w_i x_i y_i,
  \label{eq:boltzmann-sums}
\end{equation}
from which the slope and intercept are
\begin{equation}
  b = \frac{S_w S_{wxy} - S_{wx} S_{wy}}{S_w S_{wxx} - S_{wx}^2}, \quad
  a = \frac{S_{wxx} S_{wy} - S_{wx} S_{wxy}}{S_w S_{wxx} - S_{wx}^2}.
  \label{eq:boltzmann-slope}
\end{equation}
The temperature follows as $T_K = -1/(b \cdot k_B)$, and the uncertainty propagates as $\sigma_T = T_K^2 k_B \sigma_b$.

For $D$ elements with varying numbers of emission lines, direct vectorization would require ragged arrays.
We instead adopt a pad-and-mask strategy: all elements are padded to a common maximum line count $N_\mathrm{max}$, with a binary mask $m_i \in \{0, 1\}$ that zeroes the weight $w_i$ for padded entries.
The five sums in \cref{eq:boltzmann-sums} are then computed as standard inner products on arrays of shape $(D, N_\mathrm{max})$, requiring $10 N_\mathrm{max} + 7$ floating-point operations per element.
This pad-and-mask approach maps directly to GPU SIMD execution without requiring \texttt{vmap}---the element dimension $D$ is simply a batch axis over which the five dot products are computed in parallel.

\subsection{Anderson-Accelerated Saha--Boltzmann Solver}
\label{sec:anderson}

The CF-LIBS inversion requires solving a charge-balance equation self-consistently: given a temperature $T$ and trial electron density $n_e$, the Saha equation determines ionization fractions, from which a new $n_e$ follows via charge neutrality $n_e = \sum_s \sum_{z \geq 1} z \cdot n_{s,z}$.
This defines a fixed-point problem $n_e = g(n_e)$ where $g$ encapsulates the Saha--Boltzmann system.

Simple Picard iteration $n_e^{(k+1)} = g(n_e^{(k)})$ converges at a rate governed by the spectral radius of $g'(n_e^*)$, which is typically $0.3$--$0.7$ for LIBS-relevant conditions.
Anderson acceleration \citep{Walker2011} improves upon Picard iteration by maintaining a history of the $m$ most recent iterates and their residuals $r_k = g(n_e^{(k)}) - n_e^{(k)}$, then computing an optimal linear combination:
\begin{equation}
  n_e^{(k+1)} = g(n_e^{(k)}) - (\Delta G_k + \Delta R_k)\, \bm{\gamma}_k^*,
  \label{eq:anderson}
\end{equation}
where $\Delta G_k$ and $\Delta R_k$ are matrices of iterate and residual differences over the history window, and $\bm{\gamma}_k^*$ minimizes the residual norm:
\begin{equation}
  \bm{\gamma}_k^* = \arg\min_{\bm{\gamma}} \| r_k - \Delta R_k \bm{\gamma} \|^2 + \lambda \|\bm{\gamma}\|^2.
  \label{eq:anderson-ls}
\end{equation}
The Tikhonov regularization parameter $\lambda$ stabilizes the least-squares solve when the history vectors become nearly collinear.

\citet{Toth2015} proved that Anderson acceleration with depth $m$ converges r-linearly at a rate improved by a factor depending on the gain over Picard's contraction constant.
\citet{Evans2018} showed that for sufficiently smooth fixed-point maps, depth~$m$ Anderson acceleration achieves convergence comparable to GMRES($m$) applied to the linearized residual equation.

In our JAX implementation, the iteration history is maintained as a fixed-size circular buffer within a \texttt{lax.scan} loop, ensuring a static computation graph suitable for JIT compilation.
The depth $m$ is a compile-time constant (we find $m = 1$ optimal at tolerance $10^{-6}$ for typical LIBS conditions, balancing the overhead of the least-squares solve against the iteration reduction).
When $m = 0$, the algorithm reduces exactly to Picard iteration, providing a graceful fallback.

\subsection{Softmax Compositional Closure}
\label{sec:softmax}

CF-LIBS determines elemental compositions subject to the simplex constraint $\sum_{i=1}^D C_i = 1$ with $C_i > 0$.
Traditional approaches enforce this constraint via explicit normalization after each iteration, which complicates gradient-based optimization and can introduce numerical instability when compositions approach the simplex boundary.

The softmax reparameterization maps unconstrained parameters $\bm{\theta} \in \mathbb{R}^D$ to the probability simplex via
\begin{equation}
  C_i = \frac{\exp(\theta_i)}{\sum_{j=1}^D \exp(\theta_j)},
  \label{eq:softmax}
\end{equation}
implemented with the log-sum-exp trick for numerical stability: $\theta_i \leftarrow \theta_i - \max_j \theta_j$ before exponentiation.
The constraint $\sum_i C_i = 1$ is satisfied exactly by construction, and strict positivity $C_i > 0$ holds for all finite $\bm{\theta}$.

The Jacobian of the softmax map,
\begin{equation}
  J_{ij} = \frac{\partial C_i}{\partial \theta_j} = C_i(\delta_{ij} - C_j),
  \label{eq:softmax-jacobian}
\end{equation}
is rank $D - 1$ (the kernel is $\bm{1}$, reflecting the shift invariance $C(\bm{\theta}) = C(\bm{\theta} + c\bm{1})$) and has condition number $\kappa = 1$ at the uniform composition $C_i = 1/D$.
This Jacobian propagates gradients from any downstream loss function (e.g., spectral residual) back through the compositional constraint, enabling joint optimization of $T$, $n_e$, and $\bm{C}$ via L-BFGS-B or similar quasi-Newton methods.

\subsection{Batch Forward Model}
\label{sec:batch}

The forward model maps plasma parameters $(T, n_e, \bm{C})$ and a wavelength grid $\bm{\lambda}$ to a synthetic emission spectrum $S(\lambda)$ through four stages:
\begin{enumerate}
  \item \textbf{Saha ionization balance}: For each element~$s$, compute the ionization-stage fractions $f_z^{(s)}(T, n_e)$ from the Saha equation using partition functions $U_z(T)$ and ionization potential depression $\Delta\chi_\mathrm{DH}$.
  \item \textbf{Boltzmann level populations}: For each transition line~$l$ belonging to element~$s$ in stage~$z$, compute the upper-level number density $n_k^{(l)} = C_s \cdot n_e \cdot f_z^{(s)} \cdot g_k^{(l)} \exp(-E_k^{(l)} / T) / U_z^{(s)}(T)$.
  \item \textbf{Line emissivity}: Compute $\varepsilon_l = (hc / 4\pi\lambda_l) A_{ki}^{(l)} n_k^{(l)}$ in \si{\watt\per\cubic\metre\per\steradian} (with the \si{\per\cubic\centi\metre} to \si{\per\cubic\metre} conversion factor $10^6$).
  \item \textbf{Voigt broadening and spectrum assembly}: Evaluate the Voigt profile for each line on the wavelength grid via \cref{eq:voigt-broadcast}, then sum:
    \begin{equation}
      S(\lambda) = \sum_{l=1}^{N_\mathrm{lines}} \varepsilon_l \cdot V_l(\lambda) \quad [\si{\watt\per\cubic\metre\per\steradian\per\nano\metre}].
      \label{eq:spectrum}
    \end{equation}
\end{enumerate}

For batch computation over $B$ parameter sets $\{(T_b, n_{e,b}, \bm{C}_b)\}_{b=1}^B$---as required for manifold generation or Monte Carlo uncertainty propagation---we apply JAX's \texttt{vmap} transformation at a single level over the batch dimension:
\begin{equation}
  \texttt{batch\_spectrum} = \texttt{jit}(\texttt{vmap}(\texttt{single\_spectrum},\; \texttt{in\_axes}=(0, 0, 0, \text{None}, \text{None}))),
  \label{eq:vmap}
\end{equation}
where the first three axes (temperature, electron density, composition) are mapped over the batch, while the wavelength grid and atomic data are broadcast (shared) across all batch elements.

The dominant memory consumer is the profile matrix of shape $(B, N_\mathrm{wl}, N_\mathrm{lines})$.
The maximum batch size for a given GPU memory budget $M_\mathrm{GPU}$ is
\begin{equation}
  B_\mathrm{max} = \left\lfloor \frac{M_\mathrm{GPU} - M_\mathrm{overhead}}{N_\mathrm{wl} \times N_\mathrm{lines} \times d_\mathrm{bytes}} \right\rfloor,
  \label{eq:bmax}
\end{equation}
where $d_\mathrm{bytes} = 8$ for \texttt{float64} and $M_\mathrm{overhead} \approx \SI{2}{\giga\byte}$ accounts for the JAX runtime and XLA workspace.

The batch forward model achieves bit-identical results to sequential single-spectrum evaluation because \texttt{vmap} is a purely semantic transformation---it does not alter the floating-point operations within each spectrum computation, only their scheduling across GPU threads.
This property provides a strong correctness guarantee: any result obtained in batch mode is numerically identical to the same computation performed sequentially on CPU.
As reported by \citet{Kawahara2022} for ExoJAX, this \texttt{vmap}-based approach to spectral batch computation is well-suited to GPU architectures where the per-spectrum workload (dominated by the $(N_\mathrm{wl}, N_\mathrm{lines})$ profile matrix) provides sufficient arithmetic intensity to saturate the GPU's compute units.
